\section{재현 방법}\label{app:repro}
재현 패키지(ZIP)를 풀거나 \url{https://github.com/rnjswnghl/reflex-hub}를 받은 뒤 아래를 실행하면, 원자료에서 이 논문의 모든 수치와 그림이 다시 만들어진다.
자세한 내용은 패키지의 \texttt{REPRODUCE.md}에 있다.

\begin{enumerate}
\item 환경: Python 3.10 이상, \texttt{pip install -r analysis/requirements.txt}. 논문 빌드에는 XeLaTeX이 필요하다.
\item 원자료 정리: \texttt{python3 analysis/build\_dataset.py}\\
\texttt{data/raw/sessions/*.json} $\to$ \texttt{data/raw/trials.csv}, \texttt{sessions.csv}
\item 사전 등록한 분석: \texttt{python3 analysis/run\_analysis.py}\\
$\to$ \texttt{analysis/output/results.json}, 논문 수치 매크로, 그림
\item 논문 빌드: \texttt{bash scripts/build\_paper.sh}
\item 실험 페이지 직접 실행: \texttt{cd web \&\& npm install \&\& npm run dev}
\end{enumerate}

\section{AI에게 맡긴 일과 직접 판단한 일}\label{app:ai}
이 연구는 AI(Claude)와 함께 진행했다.
역할 구분과 판단 근거는 재현 패키지의 \texttt{docs/06-ai-and-me.md}에 기록했고, 요약은 다음과 같다.
AI는 가설 문장의 결함 지적, 실험 페이지와 분석 코드 작성, 문헌 후보 제시와 서지 대조, 원고 초안 작성을 맡았다.
AI가 제시한 문헌은 모두 웹에서 직접 서지를 대조했고, 대조 과정에서 원문과 맞지 않는 연결 문장 1건을 고쳤다.
분석 규칙은 데이터를 보기 전에 고정했으므로, AI도 연구자도 결과를 본 뒤 판정 기준을 바꿀 수 없었다.

\section{문서 대응표}
과제 통과 기준(T10-C01--C51)과 이 논문·패키지의 대응 위치는 \texttt{docs/CHECKLIST.md}에 있다.
